\documentclass[spanish,letterpaper,oneside]{article}

\def\documenttitle {Repensar la Escuela: Del Modelo Tradicional al Aprendizaje Profundo}
\def\documentsubtitle {Fundamentos para la enseñanza y el aprendizaje I · Cohorte SD26}
\def\documentsubject {Actividad 3 · Infografía y Análisis (Módulo 51440)}
\def\documentauthor {Mtro. Martín Jonathan de la Cruz Muñoz}
\def\coursename {Fundamentos para la enseñanza y el aprendizaje I}
\def\coursecode {MAEM-SD26}

\def\universityname {Instituto de Investigación, Innovación y Estudios de Posgrado para la Educación del Estado de Nuevo León}
\def\universityfaculty {Maestría en Aprendizaje y Enseñanza de las Matemáticas}
\def\universitydepartment {Fundamentos para la enseñanza y el aprendizaje I}
\def\universitydepartmentimage {img/iiiepe/marca-iiiepe.png}
\def\universitydepartmentimagecfg {height=1.75cm}
\def\universitylocation {Monterrey, Nuevo León}
\def\studentfullname {Mtro. Martín Jonathan de la Cruz Muñoz}
\def\studentid {00874}
\def\studentemail {martin.delacruz@campus.iiiepe.edu.mx}
\def\programperiod {2do periodo, septiembre--diciembre de 2026 (SD26)}
\def\programmodality {Modalidad presencial}
\def\programshort {MAEM}

\def\authortable {
  \begin{tabular}{ll}
    \textbf{Alumno:} & \studentfullname \\
    Matrícula: & \studentid \\
    Correo institucional: & \studentemail \\
    Unidad: & \coursename \\
    Programa: & \programshort \\
    Periodo: & \programperiod \\
    & \\
    \multicolumn{2}{l}{Fecha de realización: 19 de noviembre de 2026} \\
    \multicolumn{2}{l}{\universitylocation}
  \end{tabular}
}

\input{template}

\usepackage{helvet}
\renewcommand{\familydefault}{\sfdefault}
\usepackage{array}
\usepackage{longtable}
\usepackage{pdflscape}
\usepackage{tikz}
\usepackage{pgfplots}
\pgfplotsset{compat=1.18}
\usepackage[most]{tcolorbox}
\usetikzlibrary{arrows.meta,positioning,calc,fit,backgrounds,shapes.geometric}

\definecolor{iiiepeNavy}{HTML}{163B5C}
\definecolor{iiiepeBlue}{HTML}{1F6F9F}
\definecolor{iiiepeAqua}{HTML}{20B7A8}
\definecolor{iiiepeOrange}{HTML}{F28C28}
\definecolor{iiiepePaper}{HTML}{F6F8F7}
\definecolor{iiiepeInk}{HTML}{1E2A32}
\definecolor{iiiepePale}{HTML}{EEF5F8}
\definecolor{iiiepeMint}{HTML}{E7F6F4}
\definecolor{iiiepeWarm}{HTML}{FFF3E6}
\def\portraittitlecolor {iiiepeNavy}
\setcitestyle{authoryear,open={(},close={)}}
\makeatletter
\renewcommand{\@biblabel}[1]{}
\makeatother
\setlength{\bibhang}{1.27cm}
\setlength{\bibsep}{6pt}
\AtBeginDocument{\renewcommand{\refname}{Referencias}}

\newcommand{\insertcoverwatermark}{%
  \AddToShipoutPictureBG*{%
    \begin{tikzpicture}[remember picture,overlay]
      \node[opacity=0.04,inner sep=0pt] at (current page.center) {%
        \includegraphics[width=0.65\paperwidth]{img/iiiepe/marca-iiiepe.png}%
      };
    \end{tikzpicture}%
  }%
}

\setlength{\parskip}{5pt}\setlength{\parindent}{0pt}
\begin{document}\insertcoverwatermark\templatePortrait\templatePagecfg\templateFinalcfg\fancyfoot{}\fancyfoot[L]{\scriptsize IIIEPE · MAEM · SD26}\fancyfoot[C]{\scriptsize\thepage}\fancyfoot[R]{\scriptsize Actividad 3 · Infografía}\setstretch{1.3}\raggedbottom
\section{Introducción y postura crítica compartida}
La lectura articulada de \citet{meirieu2020escuela} y \citet{sawyer2008optimizing} en torno al eje \textit{Repensar la escuela: del modelo tradicional al aprendizaje profundo} confronta una de las inercia más persistentes de la educación media superior y superior: suponer que basta con ajustar herramientas o dispositivos digitales para mejorar el aprendizaje sin transformar la arquitectura pedagógica de la clase. Sostengo la tesis de que tanto la justicia educativa planteada por Meirieu como la comprensión conceptual fundamentada por las ciencias del aprendizaje de Sawyer convergen en un mismo imperativo: la escuela solo cumple su mandato democrático cuando pasa de la transmisión individualizada y rutinaria a comunidades dialógicas de aprendizaje profundo guiadas por mediación docente experta.

En esta Actividad 3 (módulo \texttt{51440}), presento primero la infografía solicitada en formato estructurado —con título, cuerpo visual multi-bloque, fuentes y créditos— integrando la totalidad de conceptos clave examinados en ambas lecturas y en la presentación de apoyo del curso. Posteriormente desarrollo un análisis propio de sus tensiones, su vínculo con la enseñanza de las matemáticas en contextos técnicos y de ingeniería, y sus conclusiones prácticas.

\clearpage
\begin{landscape}
\thispagestyle{empty}
\section*{Infografía · Repensar la escuela: del modelo tradicional al aprendizaje profundo}
\vspace{-4pt}
\begin{center}
\resizebox{0.97\linewidth}{!}{%
\begin{tikzpicture}[font=\scriptsize, >=Latex]
  % Marco principal apaisado: ancho total 22.8 cm, alto total 14.3 cm
  \fill[iiiepePaper] (-11.4,-7.05) rectangle (11.4,7.25);
  \draw[iiiepeNavy,very thick,rounded corners=6pt] (-11.4,-7.05) rectangle (11.4,7.25);

  % Encabezado (Titulo)
  \fill[iiiepeNavy,rounded corners=5pt] (-11.3,5.48) rectangle (11.3,7.15);
  \node[anchor=west,text width=17.4cm,align=left] at (-11.0,6.32) {%
    {\color{iiiepeAqua!45!white}\scriptsize\bfseries INFOGRAFÍA ACADÉMICA · FUNDAMENTOS PARA LA ENSEÑANZA Y EL APRENDIZAJE I (SD26)\par}%
    \vspace{2pt}%
    {\color{white}\normalsize\bfseries REPENSAR LA ESCUELA: DEL MODELO TRADICIONAL AL APRENDIZAJE PROFUNDO\par}%
    \vspace{2pt}%
    {\color{white!85}\scriptsize Diálogo crítico entre Philippe Meirieu (2020) y R. Keith Sawyer (2008): hacia una docencia comprensiva y equitativa\par}
  };
  \node[anchor=east] at (10.95,6.32) {\includegraphics[height=1.05cm]{img/iiiepe/logo-iiiepe-blanco.png}};

  % BLOQUE 1: Diagnostico compartido (izquierda superior)
  \node[draw=iiiepeNavy,fill=white,thick,rounded corners=4pt,text width=6.55cm,minimum height=4.35cm,anchor=north west,inner sep=5pt,align=left,font=\fontsize{6.8}{8.3}\selectfont] (b1) at (-11.05,5.25) {%
    {\colorbox{iiiepeNavy}{\color{white}\bfseries\ 1. DIAGNÓSTICO COMPARTIDO DE LA ESCUELA \ }}\\[4pt]
    \textbf{\color{iiiepeNavy}Philippe Meirieu (2020):} Denuncia la \textbf{brecha entre discurso y práctica}: las desigualdades escolares se naturalizan cuando se vuelve acríticamente a la «normalidad» transmisiva. La escuela pierde su sentido si reduce la educación a instrucción doméstica fragmentada o consumo individual de pantallas.\\[4pt]
    \textbf{\color{iiiepeBlue}R. Keith Sawyer (2008):} Muestra que el modelo escolar heredado de la era industrial se diseñó para transmitir hechos estandarizados, memorizar procedimientos y medir recuerdo aislado; dicho esquema resulta incompatible con la sociedad del conocimiento y bloquea la comprensión real.
  };

  % BLOQUE 2: Cambio de paradigma (centro superior)
  \node[draw=iiiepeBlue,fill=white,thick,rounded corners=4pt,text width=6.85cm,minimum height=4.35cm,anchor=north,inner sep=5pt,align=left,font=\fontsize{6.6}{8.1}\selectfont] (b2) at (0.0,5.25) {%
    {\colorbox{iiiepeBlue}{\color{white}\bfseries\ 2. CAMBIO DE PARADIGMA EDUCATIVO \ }}\\[4pt]
    \setlength{\tabcolsep}{3pt}\renewcommand{\arraystretch}{1.15}
    \begin{tabular}{p{2.95cm}|p{3.25cm}}
      \textbf{\color{iiiepeOrange!85!black}Modelo tradicional} & \textbf{\color{iiiepeNavy}Aprendizaje profundo} \\
      \hline
      Saber como datos aislados y algoritmos sin sentido & Conocimiento conceptual conectado y transferible \\
      \hline
      Estudiante receptor pasivo y evaluación punitiva & Agente que explica, conjetura y reflexiona con evidencia \\
      \hline
      Uniformidad mecánica y competencia individual & Inclusión efectiva, andamiaje y colaboración auténtica \\
      \hline
      Tecnología como simple reemplazo de pizarrón & Mediación pedagógica con propósito intelectual \\
    \end{tabular}
  };

  % BLOQUE 3: Cinco condiciones de las ciencias del aprendizaje (derecha superior)
  \node[draw=iiiepeAqua!80!black,fill=iiiepeMint,thick,rounded corners=4pt,text width=6.55cm,minimum height=4.35cm,anchor=north east,inner sep=5pt,align=left,font=\fontsize{6.8}{8.3}\selectfont] (b3) at (11.05,5.25) {%
    {\colorbox{iiiepeAqua!80!black}{\color{white}\bfseries\ 3. CINCO PRINCIPIOS DE SAWYER (2008) \ }}\\[4pt]
    \textbf{1. Conocimientos previos:} Toda enseñanza sólida parte de las ideas iniciales y concepciones alternativas del grupo.\\[2pt]
    \textbf{2. Problemas auténticos:} El razonamiento emerge ante situaciones significativas y no frente a copias mecánicas.\\[2pt]
    \textbf{3. Andamiaje progresivo:} Apoyos estructurados que se retiran gradualmente conforme crece la autonomía.\\[2pt]
    \textbf{4. Reflexión y metacognición:} Hacer visible y verbalizar por qué funciona una estrategia o dónde falló.\\[2pt]
    \textbf{5. Transferencia situada:} Aplicar principios matemáticos y científicos en escenarios nuevos.
  };

  % BLOQUE 4: Aprender juntos (izquierda inferior)
  \node[draw=iiiepeOrange!90!black,fill=iiiepeWarm,thick,rounded corners=4pt,text width=10.55cm,minimum height=4.05cm,anchor=south west,inner sep=6pt,align=left,font=\fontsize{7.0}{8.5}\selectfont] (b4) at (-11.05,-5.82) {%
    {\colorbox{iiiepeOrange!85!black}{\color{white}\bfseries\ 4. APRENDER JUNTOS: YO + NOSOTROS Y CONOCIMIENTO DISTRIBUIDO \ }}\\[4pt]
    \textbf{\color{iiiepeNavy}Lo común y lo singular (Meirieu):} La escuela existe para hacer sociedad; su núcleo no es sumar aprendizajes aislados, sino equilibrar el desarrollo singular de cada alumno con la construcción de reglas, argumentos y bienes comunes (\textit{yo + nosotros}).\\[4pt]
    \textbf{\color{iiiepeBlue}Colaboración auténtica (Sawyer):} Las investigaciones cognitivas confirman que argumentar hipótesis entre pares, confrontar representaciones distintas y distribuir responsabilidades intelectuales produce comprensiones más estables que el estudio solitario. En matemáticas, justificar ante otros transforma la fórmula en razonamiento compartido.
  };

  % BLOQUE 5: Rol docente, tecnologia, equidad y evaluacion (derecha inferior)
  \node[draw=iiiepeNavy,fill=iiiepePale,thick,rounded corners=4pt,text width=10.55cm,minimum height=4.05cm,anchor=south east,inner sep=6pt,align=left,font=\fontsize{7.0}{8.5}\selectfont] (b5) at (11.05,-5.82) {%
    {\colorbox{iiiepeNavy}{\color{white}\bfseries\ 5. DOCENTE INSUSTITUIBLE, EQUIDAD, TECNOLOGÍA Y EVALUACIÓN \ }}\\[4pt]
    \textbf{· Mediación humana irreemplazable:} Ninguna plataforma sustituye el vínculo pedagógico ni la lectura ética del aula (Meirieu). La tecnología es un medio subordinado al proyecto didáctico, nunca su reemplazo.\\[3pt]
    \textbf{· Equidad como derecho efectivo:} No basta con ofrecer acceso formal; se requieren andamiajes diferenciados para que quienes parten en desventaja alcancen comprensión profunda.\\[3pt]
    \textbf{· Evaluación auténtica:} Superar exámenes de repetición memorística mediante evidencias donde el estudiante explique procesos, relaciones y decisiones matemáticas con sentido.
  };

  % Flechas conectores en banda media limpia (sin cruzar textos)
  \draw[->,very thick,iiiepeNavy] (b1.east) -- (b2.west);
  \draw[->,very thick,iiiepeBlue] (b2.east) -- (b3.west);
  \draw[->,very thick,iiiepeOrange!85!black] (b1.south) -- (b4.north west);
  \draw[->,very thick,iiiepeAqua!80!black] (b3.south) -- (b5.north east);
  \draw[<->,very thick,iiiepeNavy] (b4.east) -- (b5.west);

  % Pie con Fuentes y Creditos
  \fill[iiiepeNavy,rounded corners=4pt] (-11.3,-6.95) rectangle (11.3,-6.02);
  \node[anchor=west,text width=14.0cm,align=left,font=\tiny\color{white}] at (-11.0,-6.48) {%
    \textbf{FUENTES (APA 7):} Meirieu, P. (2020). ``La escuela después''... ¿con la pedagogía de antes? \textit{MCEP}. · Sawyer, R. K. (2008). Optimizing learning: Implications of learning sciences research. En \textit{Innovating to learn, learning to innovate} (pp. 45--65). OECD. · Síntesis del curso PFM01 IIIEPE.
  };
  \node[anchor=east,text width=7.5cm,align=right,font=\tiny\color{white}] at (11.0,-6.48) {%
    \textbf{CRÉDITOS:} Mtro. Martín Jonathan de la Cruz Muñoz (Matrícula 00874)\\
    MAEM · IIIEPE · Cohorte SD26 · 19 de noviembre de 2026
  };
\end{tikzpicture}%
}
\end{center}
\end{landscape}
\clearpage

\section{Desarrollo analítico: convergencias, tensiones e implicaciones matemáticas}
La elaboración de la infografía anterior evidencia que \citet{meirieu2020escuela} y \citet{sawyer2008optimizing} abordan la transformación escolar desde escalas complementarias, en plena sintonía con el horizonte internacional descrito por \citet{scott2015futuro}. Mientras Meirieu interroga el \textit{para qué} político, ético y comunitario de la institución escolar —advirtiendo que regresar sin crítica a la rutina previa consolida la desigualdad y reduce la educación a consumo doméstico de contenidos—, Sawyer examina el \textit{cómo} cognitivo de la comprensión profunda, demostrando que las prácticas transmisivas de corte industrial producen conocimientos frágiles e inertes incapaces de transferirse fuera del examen \citep{sawyer2008optimizing}.

\subsection*{Análisis propio de las tensiones entre tecnología, equidad y rol docente}
Desde mi criterio como docente de matemáticas en educación media superior técnica y nivel superior en ingeniería, la convergencia más fértil entre ambos autores aparece en tres puntos críticos:
\begin{enumerate}
  \item \textbf{Tensión tecnológica y mediación pedagógica:} Tanto en la experiencia reciente de educación mediada por pantallas como en la incorporación de simuladores matemáticos, suele confundirse modernización operativa con innovación pedagógica. Sin embargo, como advierte \citet{meirieu2020escuela}, las herramientas digitales amplifican el aislamiento si cada estudiante interactúa solo frente a un tutorial que premia el acierto rápido. Solo bajo el diseño de problemas auténticos y andamiaje dialogado \citep{sawyer2008optimizing}, la tecnología apoya la visualización gráfica o el contraste experimental de hipótesis.
  \item \textbf{El aula como colectivo frente al individualismo:} En preparatoria técnica e ingeniería, fragmentar a los grupos en listas individuales de ejercicios fomenta la competencia memorística. Es decir, cuando cada alumno trabaja aislado pierde la riqueza argumentativa; por ello, recuperar la dialéctica entre lo común y lo singular (\textit{yo + nosotros}) implica diseñar espacios donde un equipo de mantenimiento o de análisis numérico confronte distintas estrategias para modelar un fenómeno, descubriendo en el diálogo por qué un procedimiento es más general o menos propenso a errores que otro \citep{eef2017mathematics}.
  \item \textbf{De la igualdad formal al derecho efectivo de aprender:} El análisis propio muestra que tratar por igual a estudiantes con puntos de partida desiguales perpetúa el rezago en matemáticas. Además, la evaluación punitiva oculta necesidades reales de aprendizaje; no obstante, garantizar el aprendizaje profundo requiere andamiajes accesibles, criterios explícitos de argumentación y evaluación formativa centrada en evidencias comprensivas en lugar de castigar la duda inicial \citep{pimienta2008evaluacion,unesco2015replantear}.
\end{enumerate}

\section{Transferencia a la práctica profesional y aula de matemáticas}
En el contexto laboral y docente de las matemáticas aplicadas, esta síntesis orienta una reorganización precisa de las secuencias didácticas. Al trabajar contenidos como razones de cambio, sistemas lineales o estadística descriptiva en bachillerato técnico e ingeniería, el enfoque tradicional presenta primero la fórmula, luego dos ejemplos resueltos en el pizarrón y finalmente veinte ejercicios repetitivos. En consecuencia, aplicar los cinco principios de \citet{sawyer2008optimizing} junto con la visión comunitaria de \citet{meirieu2020escuela} conlleva: (a) activar saberes previos mediante una situación física cercana; (b) trabajar en equipos pequeños con roles de argumentación compartidos; (c) ofrecer guías de andamiaje graduado que orienten la validación dimensional sin dar la respuesta; y (d) concluir con una puesta en común donde se exija justificar verbalmente y por escrito la validez del modelo frente a los datos.

\clearpage
\section{Conclusiones y síntesis argumentativa}
En conclusión, repensar la escuela frente al modelo tradicional exige articular dos dimensiones inseparables: la recuperación ética de la comunidad de aprendizaje propuesta por \citet{meirieu2020escuela} y el rigor cognitivo de las ciencias del aprendizaje sistematizado por \citet{sawyer2008optimizing}. Ninguna reforma técnica ni plataforma digital logra aprendizaje profundo si deja intactas la pasividad del estudiante, la fragmentación disciplinar y la evaluación meramente memorística.

Por tanto, asumir esta postura en mi práctica profesional significa situar el diseño de tareas matemáticas complejas, la colaboración deliberada y la mediación humana insustituible en el centro del trabajo diario, garantizando que el conocimiento científico y matemático opere como derecho efectivo para todo el grupo.

\clearpage
\begin{thebibliography}{99}
\bibitem[Henderson et~al.(2017)]{eef2017mathematics} Henderson, P., Hodgen, J., Foster, C., y Kuchemann, D. (2017). \textit{Improving mathematics in Key Stages 2 and 3}. Education Endowment Foundation.
\bibitem[Meirieu(2020)]{meirieu2020escuela} Meirieu, P. (2020, 18 de abril). ``La escuela después''... ¿Con la pedagogía de antes? \textit{Movimiento Cooperativo de Escuela Popular}. \url{https://www.mcep.es/la-escuela-despues-con-la-pedagogia-de-antes-philippe-meirieu/}
\bibitem[Pimienta Prieto(2008)]{pimienta2008evaluacion} Pimienta Prieto, J. H. (2008). \textit{Evaluación de los aprendizajes: Un enfoque basado en competencias}. Pearson Educación.
\bibitem[Sawyer(2008)]{sawyer2008optimizing} Sawyer, R. K. (2008). Optimizing learning: Implications of learning sciences research. En OECD, \textit{Innovating to learn, learning to innovate} (pp. 45--65). OECD Publishing. \url{https://doi.org/10.1787/9789264039958-en}
\bibitem[Scott(2015)]{scott2015futuro} Scott, C. L. (2015). \textit{El futuro del aprendizaje 1: ¿Por qué deben cambiar el contenido y los métodos de aprendizaje en el siglo XXI?} UNESCO.
\bibitem[UNESCO(2015)]{unesco2015replantear} UNESCO. (2015). \textit{Replantear la educación: ¿Hacia un bien común mundial?} Ediciones UNESCO.
\end{thebibliography}

\end{document}
